\section{Bitter Lesson：不要和模型时代最大变量背道而驰}

上一章讲产品底盘，本章进入 AI 应用创业的第一次大转向。明超平从 Moonshot 出来创业后，很快遇到 bitter lesson：当模型能力是时代最大变量时，应用公司如果本能地给产品加太多脚手架、模板和雕花，可能反而背离模型进步方向。产品还没上线就停掉，不是因为没有做事，而是意识到做错了变量。

\begin{figure}[H]
\centering
\includegraphics[width=0.96\textwidth]{bitter-lesson-pivot.png}
\caption{Bitter Lesson 转向：想雕花和加脚手架，可能背离时代最大变量。自制概念图，依据 01:05:05--01:13:18 对谈内容整理。}
\end{figure}

\begin{knowledgebox}{读图：叫停也是产品能力}
图中从 Scaffold 到 Launch?、Wrong variable、Stop 和 Rethink。这里的关键不是“什么都别做”，而是应用公司要判断哪些体验应该由产品承接，哪些能力应该让模型自然长出来。错把脚手架当壁垒，会让产品越来越重。
\end{knowledgebox}

\subsection{Noisee 的走红与骤停}

本节补上 Moonshot 阶段的一个重要转折。访谈里提到海外产品 Noisee 的走红与骤停，它让明超平看见 AI 产品窗口的两个特点：爆发可能来得很快，但产品是否真的顺着模型能力和用户场景走，需要更冷静地判断。一个产品可以在传播上得到短期正反馈，但如果底层变量不对，继续雕花只会让团队离时代主线更远。

这段经历和 bitter lesson 互相解释。AI 应用创业不是“做一个好看的 demo 然后投流”，而是要判断模型正在释放哪类能力，用户在哪个环境里能最快把能力变成作品或生产力。Noisee 给了他前沿体感，YouWare 则是重新选择了更贴近 coding 创作和发布环境的切口。

\begin{warningbox}{传播正反馈不等于产品方向正确}
AI 产品很容易因为新奇感获得传播，但传播不能替代留存、创作、复用和生态。明超平在 Moonshot 和创业早期的转折，正是在区分“火了”与“方向对了”。
\end{warningbox}

\subsection{Token 消耗速度与 per token valuation}

失眠后的顿悟，是 AI 时代的关键指标之一可能是 token 消耗速度。用户不是为了“花 token”而花 token，而是为了让智能更快转成结果。一个好的 Agent 场景，应该用更少 token 完成更大任务，或者让 token 更高效地转化为作品、网站、游戏、流程和生产力。

\begin{figure}[H]
\centering
\includegraphics[width=0.96\textwidth]{token-consumption-speed.png}
\caption{Token 消耗速度：AI 时代产品指标从点击/时长转向更高效消耗智能。自制概念图，依据 01:13:18--01:16:33 与 02:03:33--02:14:50 对谈内容整理。}
\end{figure}

\begin{knowledgebox}{读图：token 是智能消耗，不只是成本}
图中 User intent 进入 Container，再由 Agent action 消耗 Tokens，最后转成 Value。读图时不要只把 token 看作 API 成本；对应用公司来说，token 是智能被调动的单位，关键是每个 token 创造多少用户价值。
\end{knowledgebox}

\[
\mathrm{PTV}=\frac{\mathrm{User\ Value}}{\mathrm{Tokens\ Consumed}}
\]

其中，$\mathrm{PTV}$ 表示 per token valuation；$\mathrm{User\ Value}$ 是用户实际得到的作品、效率、收入或体验；$\mathrm{Tokens\ Consumed}$ 是完成任务消耗的模型 token。这个式子不是财务公式，而是产品思考工具：如果消耗很多 token 却只产生一次闲聊，价值密度低；如果少量 token 就能生成可分享、可交互、可继续编辑的网站或游戏，价值密度高。

\begin{warningbox}{不要把 token 消耗误读成越多越好}
消耗速度重要，但不能粗暴理解为烧 token。真正目标是更快、更高效地把智能转为用户可感知价值。无效长对话、反复试错和空转推理，只会提高成本，不会形成产品壁垒。
\end{warningbox}

\subsection{本章小结}

AI 应用创业的第一课是顺着模型时代的最大变量，而不是用旧产品习惯给模型套复杂壳。YouWare 的关键指标从传统点击/时长，转向智能如何被容器化、高效消耗并转成作品。

